\chapter{Quá trình phát triển và triển khai}
\label{ch:phat-trien-trien-khai}

\section{Phương pháp thực hiện}
Nhóm áp dụng mô hình phát triển phần mềm linh hoạt \textbf{Agile/Scrum} xuyên suốt quá trình thực hiện đồ án thay vì phát triển tuần tự theo kiểu Waterfall truyền thống. Công việc được chia thành các chu kỳ lặp (sprint) ngắn kéo dài \textbf{1 tuần}, mỗi sprint tập trung hoàn thiện một nhóm chức năng cụ thể thay vì dàn trải toàn bộ hệ thống cùng lúc.

Để đảm bảo quy chuẩn và tiến độ, nhóm thiết lập bảng vận hành Scrum với các yếu tố cốt lõi:
\begin{itemize}
    \item \textbf{Vai trò (Actor):} Phân định rõ Product Owner (đại diện nghiệp vụ), Scrum Master (điều phối tiến trình) và Development Team.
    \item \textbf{Sự kiện (Ceremony):} Tổ chức Sprint Planning đầu tuần, Daily Standup hàng ngày, và Sprint Review / Retrospective cuối mỗi tuần để đánh giá kết quả và cải tiến.
    \item \textbf{Tạo tác (Artifact):} Quản lý Product Backlog, Sprint Backlog và các Increment (sản phẩm tăng trưởng) có giá trị sau mỗi sprint.
    \item \textbf{Tiêu chuẩn (DoR/DoD \& Acceptance Criteria):} 
    \begin{itemize}
        \item \textit{Definition of Ready (DoR):} Yêu cầu/user story phải có mô tả rõ ràng, đầy đủ Acceptance Criteria và được thiết kế UI/UX hoàn chỉnh trước khi đưa vào Sprint.
        \item \textit{Definition of Done (DoD):} Code phải vượt qua các bài kiểm tra tự động (unit test), không có lỗi linter, được review chéo bởi ít nhất một thành viên khác và tích hợp thành công không gây lỗi trên môi trường staging.
    \end{itemize}
\end{itemize}

Thông qua việc theo dõi biểu đồ tiến độ (milestone thực tế so với kế hoạch), \textit{velocity} (tốc độ hoàn thành story points) và \textit{defect leakage} (tỷ lệ lỗi lọt ra ngoài sprint), nhóm có thể kiểm soát và sớm phát hiện rủi ro kỹ thuật. Điển hình, 3 sự cố lớn liên quan đến AI (phân loại sai giao dịch, Prompt Injection, timeout khi xử lý dữ liệu) đã được liên kết trực tiếp vào backlog và giải quyết dứt điểm trong các sprint ở giai đoạn 3 và 4, thay vì để tồn đọng đến cuối dự án.

\section{Công cụ hỗ trợ và Quy trình CI/CD}
\subsection{Công cụ hỗ trợ phát triển}
Để phối hợp làm việc nhóm và quản lý mã nguồn hiệu quả trong suốt quá trình phát triển, nhóm sử dụng các công cụ:
\begin{itemize}
    \item \textbf{Git/GitHub:} Mỗi dịch vụ (Backend, AI Service, Mobile App, Web Admin Portal) có một repository riêng, quản lý phiên bản độc lập. Trong mỗi repo, mỗi tính năng được phát triển trên một nhánh (feature branch) riêng và được xem xét qua cơ chế Pull Request.
    \item \textbf{Jira:} Quản lý backlog, phân chia công việc theo từng sprint 1 tuần và theo dõi tiến độ hoàn thành của từng thành viên.
    \item \textbf{Slack/Zalo:} Kênh trao đổi chính, dùng để họp trực tuyến định kỳ, thông báo tiến độ và xử lý vướng mắc.
    \item \textbf{Google Drive:} Lưu trữ tài liệu chung như đề cương, tài liệu đặc tả, hình ảnh thiết kế.
\end{itemize}

\subsection{Quy trình Tích hợp và Triển khai liên tục (CI/CD)}
Dựa trên cấu hình Jenkins thực tế của hệ thống, quy trình CI/CD được thiết lập tự động hóa nhằm giảm thiểu sai sót:
\begin{itemize}
    \item \textbf{CI Pipeline (Kiểm tra và Build):} Khi có Pull Request mới, Jenkins tự động \textit{checkout} mã nguồn, cài đặt thư viện (\texttt{npm install}), chạy quá trình build (\texttt{npm run build}) để phát hiện lỗi biên dịch. Sau đó, mã nguồn được đóng gói thành \textbf{Docker Image} với tag tương ứng. (Theo cấu trúc chuẩn, quy trình bao gồm cả các bước \textit{lint}, \textit{unit test}, \textit{build} và \textit{security scan}).
    \item \textbf{CD Pipeline (Triển khai - Deploy):} Khi mã nguồn được merge vào nhánh chính (\texttt{staging} hoặc \texttt{main}), pipeline CD sẽ đồng bộ mã nguồn, yêu cầu xác nhận thủ công (Manual Approval) từ người quản lý trước khi triển khai. Sau khi xác nhận, hệ thống nạp các biến môi trường (\texttt{.env.docker}), tắt các container hiện tại và khởi tạo container mới thông qua \texttt{docker compose build} và \texttt{docker compose up -d}. Quá trình deploy staging thành công sẽ là bước đệm trước khi đưa lên production.
\end{itemize}

\textbf{Tiêu chí Rollback (Phục hồi phiên bản):} Để đảm bảo tính sẵn sàng, hệ thống quy định các tiêu chí kích hoạt rollback như sau:
\begin{itemize}
    \item \textit{Smoke test thất bại:} Ứng dụng không phản hồi hoặc container bị crash liên tục sau khi khởi động.
    \item \textit{Lỗi hệ thống tăng đột biến:} Tỷ lệ lỗi (error rate) giám sát được trên production tăng bất thường ngay sau khi triển khai phiên bản mới.
    \item Khi thỏa mãn tiêu chí, quản trị viên sẽ từ chối release (hoặc kích hoạt job rollback) để hệ thống tự động kéo lại tag Docker Image của phiên bản ổn định trước đó và khởi động lại dịch vụ.
\end{itemize}

\section{Tổ chức nhóm và phân công chuyên môn}
Nhóm gồm 6 thành viên, được chia thành 3 cặp phụ trách song song 3 mảng chức năng chính của hệ thống, mỗi cặp đảm nhận trọn vẹn cả phần giao diện (Frontend) lẫn xử lý nghiệp vụ phía máy chủ (Backend) cho mảng của mình nhằm đảm bảo tính liền mạch khi tích hợp. 

\begin{table}[H]
    \centering
    \begin{tabular}{|L{5cm}|L{9cm}|}
        \hline
        \textbf{Thành viên phụ trách} & \textbf{Mảng chức năng / Phạm vi công việc} \\
        \hline
        Võ Hoàng Nguyên, Nguyễn Duy Lâm & Mô-đun Tài chính 6 lọ (6 Jars) -- Frontend Mobile và Backend API. \\ \hline
        Nguyễn Thành Nhật, Huỳnh Công Minh & Hệ thống thông báo (Notification) và các luồng thông tin người dùng cơ bản -- Frontend Mobile và Backend API. \\ \hline
        Nguyễn Vĩnh Lương, Huỳnh Tấn Lộc & Quản lý học bổng và trang quản trị Webadmin học bổng -- Frontend (Mobile/Webadmin) và Backend API. Huỳnh Tấn Lộc phụ trách thêm DevOps. \\
        \hline
    \end{tabular}

    \caption{Phân công chuyên môn theo mảng chức năng}
    \label{tab:phan-cong}
\end{table}

Nhằm khắc phục tình trạng \textit{knowledge silo} (chia cắt kiến thức) do phân công theo cặp và giảm thiểu rủi ro khi có thành viên vắng mặt, nhóm áp dụng:
\begin{itemize}
    \item \textbf{Ma trận RACI \& Code Ownership:} Định nghĩa rõ ràng ai chịu trách nhiệm chính (Responsible) và ai phê duyệt (Accountable) cho từng tính năng, không để chồng chéo quyền hạn.
    \item \textbf{Review chéo (Cross-review):} Các Pull Request của một cặp bắt buộc phải được review và approve bởi ít nhất một thành viên từ cặp khác. Quá trình này giúp phát hiện lỗi khách quan và phân tán hiểu biết về code.
    \item \textbf{Knowledge Sharing Sessions (Phiên chia sẻ kiến thức):} Được tổ chức định kỳ sau mỗi sprint để chia sẻ về kiến trúc hệ thống, các luồng tích hợp API phức tạp và đặc biệt là kinh nghiệm xử lý AI, đảm bảo mọi người đều có khả năng \textit{backup} lẫn nhau khi cần thiết.
\end{itemize}

\section{Các giai đoạn phát triển}
Đồ án được thực hiện trong khoảng thời gian từ ngày 24/08/2025 đến ngày 21/07/2026, chia thành 5 giai đoạn chính như sau:
\begin{itemize}
    \item \textbf{Giai đoạn 1 -- Khảo sát và thiết kế hệ thống (24/08/2025 -- 09/2025):} Khảo sát thực trạng các ứng dụng quản lý tài chính và kết nối học bổng hiện có, xác định phạm vi đề tài, thiết kế kiến trúc tổng thể và cơ sở dữ liệu, hoàn thiện đề cương chi tiết.
    \item \textbf{Giai đoạn 2 -- Xây dựng nền tảng cốt lõi (10/2025 -- 01/2026):} Triển khai song song ba mảng chức năng cốt lõi theo phân công tại Bảng~\ref{tab:phan-cong}: mô-đun Tài chính 6 lọ, hệ thống thông báo và luồng người dùng cơ bản, quản lý học bổng và Webadmin. Mỗi mảng được phát triển theo các sprint ngắn, review chéo qua Pull Request trên GitHub.
    \item \textbf{Giai đoạn 3 -- Tích hợp trí tuệ nhân tạo (02/2026 -- 04/2026):} Bắt đầu xây dựng AI Service (FastAPI kết hợp LangChain và Gemini API), tích hợp trợ lý ảo AI Chat, bộ phân loại giao dịch tự động và cơ chế gợi ý học bổng vào các mảng chức năng đã hoàn thiện ở Giai đoạn 2.
    \item \textbf{Giai đoạn 4 -- Hoàn thiện, tối ưu và triển khai (05/2026 -- 06/2026):} Tối ưu hiệu năng và bảo mật (chuẩn hóa Unicode tiếng Việt, chống Prompt Injection), kiểm thử toàn diện hệ thống, cấu hình build và deploy ứng dụng lên môi trường thật.
    \item \textbf{Giai đoạn 5 -- Hoàn thiện báo cáo và bảo vệ (07/2026 -- 21/07/2026):} Tổng hợp kết quả đạt được, hoàn chỉnh báo cáo khóa luận, chuẩn bị slide và demo phục vụ báo cáo trước Hội đồng.
\end{itemize}

\section{Các thách thức và giải pháp thực hiện}
Trong suốt quá trình phát triển và tích hợp hệ thống Student360, bên cạnh các vấn đề về mặt kỹ thuật, nhóm phát triển còn đối mặt với các khó khăn thách thức lớn liên quan đến yếu tố con người và sự phối hợp giữa các nhóm. Dưới đây là các thách thức cốt lõi và các giải pháp thực tế được áp dụng:

\subsection{Thách thức ở tầng trình bày di động (Mobile App)}
\begin{itemize}
    \item \textbf{Thách thức 1 -- Truyền phản hồi AI theo thời gian thực:} Do cơ chế tải dữ liệu tuần tự trên nền di động không ổn định như trên trình duyệt, việc chờ nhận đủ toàn bộ phản hồi từ mô hình ngôn ngữ lớn (Gemini AI) trước khi hiển thị tạo ra độ trễ lớn, ảnh hưởng tiêu cực đến trải nghiệm người dùng (UX).
    \item \textit{Giải pháp:} Ứng dụng cơ chế đọc dữ liệu tăng dần theo từng phần (Streaming / Server-Sent Events - SSE). Mobile App đọc dữ liệu dạng chunk và cập nhật giao diện tức thời ngay khi nhận được ký tự đầu tiên, giúp giảm cảm giác chờ đợi của sinh viên.
    \item \textbf{Thách thức 2 -- Đồng bộ dữ liệu và tính nhất quán giữa các màn hình:} Một giao dịch chi tiêu mới của sinh viên ảnh hưởng đồng thời đến nhiều giao diện độc lập (số dư các hũ 6 lọ, lịch sử giao dịch gần đây, biểu đồ báo cáo tài chính và các cảnh báo chi tiêu). Việc tự cộng trừ số liệu thủ công ở các màn hình khác nhau dễ dẫn đến sai lệch dữ liệu.
    \item \textit{Giải pháp:} Sử dụng cơ chế Query Invalidation của thư viện React Query. Frontend tự động đánh dấu toàn bộ các nhóm dữ liệu tài chính liên quan là đã lỗi thời ngay sau khi Backend ghi nhận giao dịch thành công, buộc hệ thống tự động tải lại dữ liệu ngầm và cập nhật UI đồng bộ mà không cần tải lại toàn bộ ứng dụng.
\end{itemize}

\subsection{Thách thức ở tầng AI và Bảo mật (AI Service \& Backend)}
\begin{itemize}
    \item \textbf{Thách thức 3 -- Độ trễ lớn khi gọi LLM và xử lý tiếng Việt:} Thời gian phản hồi trung bình của chatbot AI khoảng 4 giây, có thể lên tới hơn 10 giây đối với câu hỏi phức tạp. Đồng thời, việc sinh viên nhập từ viết tắt, tiếng lóng tiếng Việt hoặc lỗi Unicode không đồng nhất (NFC/NFD) làm giảm độ chính xác của bộ phân loại giao dịch.
    \item \textit{Giải pháp:} Áp dụng bộ lọc chuẩn hóa Unicode NFC ở Backend trước khi gửi yêu cầu sang AI Service. Đồng thời, thiết lập cơ chế phân loại hỗn hợp (Hybrid Classification Pipeline): kiểm tra danh mục từ khóa cá nhân hóa trong cơ sở dữ liệu trước với độ trễ 0ms; chỉ khi không khớp mới gọi LLM để suy luận ngữ nghĩa.
    \item \textbf{Thách thức 4 -- Ngăn chặn tấn công Prompt Injection và rò rỉ dữ liệu:} Người dùng có thể cố tình nhập câu lệnh tiêm nhiễm để phá vỡ ngữ cảnh an toàn của trợ lý AI, hoặc giả mạo ID người dùng (\texttt{user\_id}) trong request API để truy cập thông tin tài chính của sinh viên khác.
    \item \textit{Giải pháp:} Xây dựng chính sách bảo mật hai tầng (Policy Filter) trên AI Service để phát hiện và chặn đứng các câu lệnh tiêm nhiễm độc hại. Ở phía Backend, hệ thống tự động ghi đè tham số \texttt{user\_id} của mọi yêu cầu bằng thông tin lấy từ token JWT đã được xác thực, ngăn chặn hoàn toàn nguy cơ client tự truyền ID tự do.
\end{itemize}

\subsection{Thách thức về yếu tố con người và phối hợp liên nhóm}
\begin{itemize}
    \item \textbf{Thách thức 5 -- Rào cản giao tiếp và nguy cơ đứt gãy thông tin trong tập thể lớn:} 
    Với quy mô dự án lên tới 14 thành viên chia làm 3 nhóm độc lập, luồng thông tin giao tiếp phình to theo cấp số nhân. Nhóm phát triển phân hệ Tài chính (6 thành viên) liên tục đối mặt với việc "nhiễu" thông tin trên các kênh trao đổi chung. Các quyết định kỹ thuật hoặc thay đổi về mặt nghiệp vụ thường bị trôi hoặc hiểu sai lệch giữa các nhóm nếu không được truyền đạt rõ ràng, dẫn đến tình trạng không đồng nhất về hiểu biết hệ thống (Information Asymmetry).
    
    \item \textbf{Thách thức 6 -- Khó khăn trong việc thống nhất kiến trúc và các mô-đun dùng chung (Shared Services):} 
    Do hệ thống chia làm 3 phân hệ hoạt động trên cùng một nền tảng cơ sở dữ liệu và hạ tầng, các chức năng dùng chung như xác thực người dùng (Authentication/Authorization), hệ thống thông báo đẩy (Notification), cài đặt ứng dụng (Setup), và cấu hình hệ thống (Configuration) trở thành các điểm thắt cổ chai (bottlenecks). Việc thống nhất kiểu dữ liệu, cấu trúc DB và đặc tả API của các phần chung này đòi hỏi sự thương thảo kéo dài giữa cả 3 nhóm. Bất kỳ sự tự ý điều chỉnh nào ở một nhóm cũng có thể trực tiếp làm đổ vỡ (break) chức năng của hai nhóm còn lại.
    
    \item \textbf{Thách thức 7 -- Sự phụ thuộc chéo giữa các tính năng (Feature Dependency) và trễ tiến độ dây chuyền:} 
    Phân hệ Tài chính có sự phụ thuộc rất lớn vào các phân hệ khác (ví dụ: cần luồng đăng nhập thành công của phân hệ lõi để lấy thông tin sinh viên, và cần hệ thống thông báo để gửi cảnh báo chi tiêu). Khi các nhóm có tốc độ phát triển không đồng đều hoặc một nhóm gặp sự cố kỹ thuật dẫn đến chậm trễ tiến độ (delay), toàn bộ chuỗi tích hợp của phân hệ Tài chính cũng bị đình trệ theo. Việc quản lý sự phụ thuộc này trong mô hình phát triển song song cực kỳ phức tạp.
    
    \item \textbf{Thách thức 8 -- Xung đột mã nguồn và bất đồng bộ trong quy trình triển khai (Git Conflict \& Release Management):} 
    Với 14 lập trình viên cùng đẩy mã nguồn lên các repository cốt lõi của hệ thống (như Backend NestJS và Frontend Mobile), tần suất xảy ra xung đột code (code conflict) rất cao. Việc thiếu thống nhất trong quy chuẩn viết code (coding convention), quy trình đặt tên nhánh (Git branch naming) hay thời điểm merge code lên các nhánh \texttt{staging}/\texttt{main} gây khó khăn lớn cho việc kiểm soát phiên bản và triển khai tự động (CI/CD).

    \item \textit{Giải pháp thực hiện:} Để giải quyết triệt để các thách thức về mặt con người và quy trình phối hợp trên, các nhóm đã cùng thống nhất và áp dụng đồng bộ các giải pháp sau:
    \begin{itemize}
        \item \textit{Chuẩn hóa kênh và tần suất giao tiếp:} Thiết lập các kênh thảo luận chuyên biệt cho từng chủ đề nghiệp vụ trên Slack/Zalo. Định kỳ tổ chức các cuộc họp giao ban liên nhóm (Cross-team Sync Meeting) hàng tuần để đại diện các nhóm báo cáo tiến độ, chỉ ra các điểm phụ thuộc chéo và cùng đưa ra quyết định chung bằng văn bản.
        \item \textit{Thiết kế và đóng băng đặc tả API trước khi lập trình:} Sử dụng công cụ Swagger/OpenAPI để định nghĩa chi tiết các API endpoints dùng chung. Bất kỳ sự thay đổi nào về cấu trúc request/response phải được đăng ký qua một biểu mẫu yêu cầu thay đổi (RFC - Request for Comments) và chỉ được áp dụng khi có sự đồng thuận từ đại diện của cả 3 nhóm.
        \item \textit{Quy hoạch rõ ranh giới trách nhiệm (RACI Matrix):} Phân định rõ ràng nhóm nào chịu trách nhiệm phát triển chính và nhóm nào chịu trách nhiệm tích hợp đối với từng API hoặc màn hình dùng chung, tránh việc đùn đẩy trách nhiệm hoặc làm trùng lặp công việc.
        \item \textit{Thống nhất quy trình Git Flow và rà soát mã nguồn liên nhóm (Cross-review):} Áp dụng quy chuẩn viết mã chung (Linter/Prettier) trên toàn bộ dự án. Các Pull Request liên quan đến các mô-đun dùng chung bắt buộc phải nhận được sự kiểm duyệt (approve) của ít nhất một thành viên đại diện cho mỗi nhóm trước khi được phép merge.
        \item \textit{Tận dụng tối đa môi trường Staging chung:} Mọi phiên bản mã nguồn trước khi đưa lên production đều phải được tích hợp và kiểm thử tự động trên môi trường Staging độc lập. Điều này giúp cô lập các lỗi tích hợp sớm và đảm bảo sự cố ở một phân hệ không làm gián đoạn công việc của các phân hệ khác.
    \end{itemize}
\end{itemize}
